\chapter{Vehicle Purchase Strategy}

Buying a family car is planned for 2033. Unlike a house, a vehicle depreciates from the moment it leaves the showroom---and continues to lose value every year regardless of how well it is maintained. Getting the decision right means not just picking the right vehicle, but choosing the right method of acquiring it and understanding what it truly costs to own per month. This phase compares all three acquisition options, selects the best fit for my circumstances, and calculates total monthly ownership cost.

\section{Comparison of Acquisition Options}

\begin{table}[h!]
\caption{Comparison of Vehicle Acquisition Options}
\label{tab:vehicle_options}
\begin{tabularx}{\linewidth}{@{}lLLL@{}}
\toprule
\textbf{Factor} & \textbf{Lease} & \textbf{New Purchase} & \textbf{Used Purchase} \\
\midrule
Upfront Cost   & Low or none & High down payment & Lower down payment \\
Monthly Cost   & Lower fixed payment & Higher loan repayment & Lower loan repayment \\
Ownership      & No ownership at end & Full ownership & Full ownership \\
Depreciation Risk & Borne by lessor & High: 30--40\% Year 1 & Low: already depreciated \\
Maintenance    & Often included & Owner's full responsibility & Owner's full responsibility \\
Flexibility    & Mileage caps; no modifications & Keep, modify, or sell freely & Keep, modify, or sell freely \\
Best For       & Short-term, business use & Long-term family ownership & Budget-conscious buyers \\
\bottomrule
\end{tabularx}
\end{table}

\textbf{Why not lease?} Leasing ends without ownership---every payment builds equity in the vehicle, then hands it back at contract end. For business users who need a new vehicle every two to three years and can deduct lease costs, this makes sense. For a young family planning to keep a car for seven to ten years in Hanoi, it does not. The cumulative lease payments over a decade would likely exceed the outright purchase cost, with nothing to show for it at the end.

\textbf{Why not used?} A 3-year-old Honda Civic of equivalent specification typically trades at 650,000,000 to 700,000,000~VND in Vietnam's secondary market---a 22 to 28\% discount to the 900,000,000~VND new price. That is a real saving. However, two constraints make it the wrong choice specifically in this plan: first, used hybrid models are rare in Vietnam---the Honda Civic e:HEV is not yet widely traded second-hand, which limits selection and raises overpayment risk; second, any used vehicle bought in 2033 is at least three years old with an expired factory warranty, inheriting mechanical risk at exactly the life stage when reliability matters most.

\textbf{Selected: New Purchase.} Full manufacturer warranty, complete service history from day one, and access to the latest hybrid or EV efficiency technology. At BCG Consultant income levels in 2033, the premium over a used purchase is manageable and the peace-of-mind value is high.

\section{Vehicle Specifications and Financing}

Target: Honda Civic e:HEV (Hybrid 5-seater sedan) or VinFast VF7 (Electric SUV) at approximately 900,000,000~VND. Financing structure: 300,000,000~VND down payment (33\%), 600,000,000~VND bank loan at 8.0\% per annum over 60 months. Monthly loan repayment: approximately \textbf{12,166,000~VND}. At a projected income of 120,000,000 VND per month as BCG Consultant in 2033, this is 10.1\% of gross income---well within the 15\% vehicle debt ceiling. The 33\% down payment is deliberately above the typical 20 to 25\% bank minimum: a larger equity stake from day one reduces the risk of being underwater on the loan if resale value drops faster than expected in the first two years.

\section{Total Cost of Ownership}

\begin{table}[h!]
\caption{Monthly Vehicle Running Costs}
\label{tab:vehicle_tco}
\begin{tabularx}{\linewidth}{@{}lrL@{}}
\toprule
\textbf{Cost Item} & \textbf{Monthly (VND)} & \textbf{Notes} \\
\midrule
Fuel / Electricity      & 1,500,000 & Hybrid/EV saves 500,000--1,000,000 vs.\ petrol equivalent \\
Maintenance / Battery   & 1,200,000 & Scheduled service or VinFast battery lease fee \\
Parking \& Tolls        & 1,300,000 & Monthly parking in Hanoi + urban road tolls \\
Insurance \& Registration & 1,000,000 & Civil liability + optional comprehensive (amortized) \\
\midrule
\textbf{Total Running Costs} & \textbf{5,000,000} & \\
\textbf{Total Monthly (Loan + Running)} & \textbf{17,166,000} & 14.3\% of projected 120M income \\
\bottomrule
\end{tabularx}
\end{table}

The 14.3\% total vehicle cost as a share of income sits just below the standard 15\% guideline for vehicle debt. The hybrid or electric engine choice reduces the fuel line significantly compared to a conventional petrol sedan of similar size---a petrol alternative in Hanoi traffic would typically cost 2,000,000 to 2,500,000~VND/month in fuel, making the hybrid or EV option approximately 500,000 to 1,000,000~VND cheaper to run per month over a five-year ownership period.
